% !TeX root = ../papers/nc-p3.tex
\section{Walsh商の正規化と特殊な指定なし同型の有限復元}
\label{sec:walsh-spectral-fibers}
\begin{theorem}\label{thm:walsh-spectral-fibers}
定理\ref{thm:walsh-presentation-space}の平方行列の二重被覆の $S_3$ 商を $Y$ とする。
$Y$ は、Walsh表示の射影不変式商における閉像の正規化である。
標準正則像の逆像はちょうど $\operatorname{rank} N=2$ である。
全特殊点を含め、一つの窓の三つの切断の特異値から、対応し得る正規化点を有限リストに復元できる。
このリストに対する軌道閉包の線型比較で、指定を忘れた同型類の全ファイバーと最小周期が決まる。
同じ手順で任意の四次テンソルについて正則Walsh表示の有無も判定できる。
\end{theorem}
\begin{proof}[全ファイバーの有限性]
まず複素数体を使う。Walsh基底を正規直交としたテンソルは全一位置のGram行列が
スカラーで、ノルム最小ベクトルである。従って全ての非零表示が半安定であり、
同じ射影商の点を与えることとユニタリな四基底変更・共通倍率で移ることが同値である
（Kempf--Nessの最小ベクトル定理）。特に三切断の特異値を共通規格化して比較できる。
$\sum|m_{be}|^2=1$、すなわちテンソルのノルムを2に取る。

二つの隣接切断の各非零色のブロックは二次元行列表現の二重である。
四つの符号付き置換はユニタリな有限群の作用なので、直和分解もユニタリに取れる。
従って各ブロックには同じ正の特異値が二つあり、その平方は行または列の二乗和の2倍である。
零の行・列も許して三つの行和 $R_i=\sum_j|m_{ij}|^2$ と三つの列和 $C_i$ を
得る。添字順は有限通り全て取る。反対切断は各色でユニタリな対角符号と巡回行列の積である。
四つの単独枠と、三対の二重枠に十六個の特異値平方を割り当てる。単独枠は
\[
 L_\chi=\left|\sum_{i=0}^2(-1)^{i\cdot\chi}m_{ii}\right|^2
\]
であり、対 $ij=01,12,20$ の二つの値は
$L_{ij}^\pm=|m_{ij}\pm m_{ji}|^2$ でそれぞれ二度現れる。重複値、零も重複度を保って割り当てる。

$S_{ij}=(L_{ij}^++L_{ij}^-)/2$ とすると
\[
 |m_{ii}|^2=\tfrac12(R_i+C_i-\sum_{j\ne i}S_{ij}),\qquad
 \Re(m_{ii}\overline{m_{jj}})=\tfrac18\sum_\chi(-1)^{(i+j)\cdot\chi}L_\chi
 \quad(i\ne j).
\]
よって三つの複素対角成分を平面実ベクトルとしたGram行列が定まる。
その実現は非負で階数2以下のときに限り、共通複素位相を除けば共役を含め高々二通りである。
$d_i=m_{ii}^2$ とし、巡回順の対 $ij$ に対し $U_{ij}=d_k-d_i-d_j$（$k$ は残る添字）、
$J=N_{01}-N_{10}$ と置く。行列の和の条件から
\[
 N_{ij}=\tfrac12(U_{ij}+J),\qquad N_{ji}=\tfrac12(U_{ij}-J).
\]
三対の積 $L_{ij}^+ L_{ij}^-$ は全て $|J|^2$ である。この一致を確認する。
$S_{ij}=(|U_{ij}+J|+|U_{ij}-J|)/2$ なので
\[
 \Re(U_{ij}\bar J)^2=S_{ij}^2\big(|U_{ij}|^2+|J|^2-S_{ij}^2\big).
\]
$U_{ij}\ne0$ が一つでもあれば、その値と $|J|$ が決まっているため
$J$ は高々四点である。二乗で生じた余分な候補は実際の特異値比較で捨てる。
全て零なら $d_0=d_1=d_2=0$ で、自由な共通位相を使って $J=|J|$ に取れる。
こうして候補 $N$ は必ず有限個となる。各 $N$ では $P^2=\prod_{i\ne j}N_{ij}$ の根を取り、
$S_3$ と共通倍率を除く。前節の符号判定は零成分がある場合にも使える。
非対角成分が零なら未使用の符号で積を調整でき、全平方根表示が同値となる。
\path{computations/verify_walsh_spectral_fibers.py} は逆公式と重複固有値の割当てを実装し、
零対角と実でない $J$ の両方を検算する。
\end{proof}

\begin{proof}[正規化と最終比較]
九平方根の射影表示の座標環で許される対角符号を割ると、不変単項式は
全指数が偶数のものと、六非対角の指数だけが奇数のものに尽きる。
有限群の不変部分を取る操作は標数零で完全であり、
平方の五つの線型関係を課した後も被覆は $P^2=\prod_{i\ne j}N_{ij}$ である。
この二重被覆は相異なる六平面に沿う平方自由な分岐を持ち、正規である
（超曲面であり、余次元1では正則）。その有限 $S_3$ 商も正規である。
全非零表示で半安定なので、射影商への写像は $Y$ 上の射として定まる。
先の有限復元により全幾何ファイバーが有限で、射影的準有限射ゆえ有限である。
一般点では定理\ref{thm:walsh-presentation-space}の同型判定で一対一なので標数零で双有理であり、
正規化である。階数2なら三行三列は全て非零であり、正則像の記述は
定理\ref{thm:walsh-normal-quadrics}から従う。これはWalsh表示を持つ像の中での主張である。

有限リストの真のファイバーを選ぶため、任意の二テンソル $U,V$ に対して使える
基底未知数を消去した線型比較を記す。$d$ 次の対称冪で $U^d$ から始め、
四位置の行列単位をLie代数として反復して作用させ、線型包が増えなくなるまで取る。
得る空間 $L_d(U)$ は連結群 $\mathrm{GL}_4^4$ の軌道の $d$ 乗の線型包である。
標数零であるからLie代数不変部分空間はこの連結群で不変となる。
その零化空間は射影軌道閉包の $d$ 次の斉次方程式に等しい。
従って低次数から二空間を比較して違えば直ちに除ける。
終了を保証するには両軌道閉包の斉次イデアルをパラメータ表示から消去して比較してよく、
線型演算のみを使う場合は一様な有限上界 $d=4^{64}$ も使える。

実際、全行列による軌道写像は64変数で次数4である。写像のグラフに64個の一般線型条件を
課すとBézoutの次数上界は $4^{64}$ であり、線型射影した閉像の次数もこれ以下である。
可逆行列は稠密なのでその射影閉像が所望の軌道閉包である。
次数 $D$ の射影既約多様体の外の点を分離するには高々 $D$ 次の方程式で足りる。
次元を $r$ とすると、その点を通る一般の余次元 $r+1$ の線型空間は多様体と交わらず、
その中の超平面から $\mathbb P^{r+1}$ へ射影すれば、像の超曲面の次数は高々 $D$ で
対象点の像を含まない。方程式を引き戻せばよい。次数を上げるにはその点で非零の線型式を掛ける。
ゆえに上界で二つの $L_d$ が等しければ閉包が等しい。
各軌道はその既約閉包に稠密な開集合なので、閉包が等しいことと同型は同値である。
この最悪上界は実際の計算量の小ささを主張するものではない。

入力が任意のテンソルなら、まず四位置のGram行列がスカラーとなる代表を求める。
複素行列の実・虚部について、この条件と行列の可逆性は実代数的であり、
実閉体上の消去で解の有無と一つの代表を決定できる。代表が存在しなければWalsh軌道ではない。
存在すれば上の三つのスペクトルから有限個の $N,P$ を作り、実際の表示テンソルと
$L_d$ または軌道イデアルを比較する。隣接切断の必要な重複度も最初に照合する。
標数零の他の代数閉体上でも、入力の有限生成係数体を複素数へ埋め込んで比較できる。
同型の有無は有限型方程式の幾何学的可解性であり、選んだ埋込みや共役の選択に依存しない。
正則な分岐内の最小周期は、実際のファイバーが
$(N^{\mathsf t},-P)$ を含めば1、それ以外は2である。
\end{proof}
